% NeurIPS Paper Checklist — \input this at the end of the NeurIPS-formatted paper
% (after \bibliography, before \end{document}). Answers are honest to the current
% state; update the [TODO] items before submission.
\newpage
\section*{NeurIPS Paper Checklist}

\begin{enumerate}
\item {\bf Claims} \\
Question: Do the abstract and introduction accurately reflect the paper's contributions and scope? \\
Answer: \answerYes{} \\
Justification: The abstract states the contribution as the instrument and its measurements, explicitly frames the criticality/edge-of-chaos phenomena as decades-old (measured, not discovered), and flags that the capacity effect is masked-specific and does not replicate on autoregressive models. Every headline claim is hedged to its statistics: the capacity ordering is presented as a suggestive two-seed gap (tiny$\ll$\{mini,base\}), not an established scaling axis; the front velocity as an apparatus-kinematic law; the AR capacity trend as a null. A withdrawn kinematics$\perp$stability decomposition is disclosed rather than retained.

\item {\bf Limitations} \\
Question: Does the paper discuss the limitations of the work? \\
Answer: \answerYes{} \\
Justification: The Limitations section states: globally-inconsistent MLM conditionals (a well-defined dynamical system but not a sampler of any joint); small models ($\le$110M masked, $\le$1B AR) and small rings ($N\le384$); proxy census (WikiText, not the models' training data); the AR capacity trend rests on few models; the special-token scheme is a first-class apparatus; temperature scales are not comparable across constructions.

\item {\bf Theory assumptions and proofs} \\
Question: For each theoretical result, does the paper provide the full set of assumptions and a complete (and correct) proof? \\
Answer: \answerNA{} \\
Justification: The paper contains no formal theorems; it is empirical/measurement work. Connections to the Lieb--Robinson bound and the Bagnoli--Rechtman--Ruffo Lyapunov--velocity relation are cited, not re-derived.

\item {\bf Experimental result reproducibility} \\
Question: Does the paper fully disclose all the information needed to reproduce the main experimental results? \\
Answer: \answerYes{} \\
Justification: All code, model checkpoints (toy), per-run result JSON/NPZ, and figures are in the public repository; the Reproducibility appendix gives models, grids, seeds, batch/ring sizes, the diversity-floor and CRN definitions, and the statistical tests. A harness regression test asserts the CRN null coupling is exactly zero.

\item {\bf Open access to data and code} \\
Question: Does the paper provide open access to the data and code, with sufficient instructions? \\
Answer: \answerYes{} \\
Justification: Repository with a README giving per-phase reproduction commands and M1 runtime tables; pinned dependency versions; scripts are idempotent and resumable. [TODO: make the repository public and add the URL at camera-ready.]

\item {\bf Experimental setting/details} \\
Question: Does the paper specify all the training and test details? \\
Answer: \answerYes{} \\
Justification: Reproducibility appendix specifies optimizer/steps for the toy model, the exact $(r,T,N,B,\text{sweeps},\text{seeds})$ grids per experiment, the reference corpora, and the metric definitions.

\item {\bf Experiment statistical significance} \\
Question: Does the paper report error bars / statistical significance? \\
Answer: \answerYes{} \\
Justification: The paper reports seed error bars and trend tests and, following an internal adversarial audit, explicitly \emph{retracts} a pseudoreplicated signed-rank statistic on the capacity claim, restating it as a suggestive two-seed gap (tiny$\ll$\{mini,base\}; base-vs-mini null $p=0.21$; AR size trend Spearman $\rho=0.17$, $p=0.29$). Marginal and null results are labeled as such. [TODO: the tracked pre-camera-ready rebuild raises the headline capacity/damping-length claims to $\ge15$ \emph{independent} seeds with a seed-level (hierarchical/bootstrap) test and a multiplicity correction; see repository issues \#1--\#3, \#9.]

\item {\bf Experiments compute resources} \\
Question: Does the paper provide sufficient information on compute resources? \\
Answer: \answerYes{} \\
Justification: All experiments run on a single Apple M1 Pro (16\,GB); per-phase wall-clock runtimes are tabulated. The toy model uses CPU JAX; real models use fp16 on the MPS backend. Total compute is a few laptop-days.

\item {\bf Code of ethics} \\
Question: Does the research conform with the NeurIPS Code of Ethics? \\
Answer: \answerYes{} \\
Justification: Measurement/interpretability work on public pretrained models and synthetic data; no human subjects, no sensitive data, no new capability that raises misuse concerns.

\item {\bf Broader impacts} \\
Question: Does the paper discuss potential positive and negative societal impacts? \\
Answer: \answerYes{} \\
Justification: The instrument is a black-box interpretability tool (positive: measurement of trained-LM token dynamics without weight access, applicable to closed models). It creates no new generative capability; negative societal impact is minimal and indirect. [TODO: one broader-impacts paragraph.]

\item {\bf Safeguards} \\
Question: Does the paper describe safeguards for responsible release of high-risk assets? \\
Answer: \answerNA{} \\
Justification: No high-risk models or datasets are released; only measurement code and analysis results on public models.

\item {\bf Licenses for existing assets} \\
Question: Are existing assets (code, data, models) properly credited and licenses respected? \\
Answer: \answerYes{} \\
Justification: Pretrained models (BERT, Pythia), datasets (tinyshakespeare, WikiText-103), and libraries (JAX, PyTorch/Transformers) are cited; usage is within their licenses. [TODO: list exact licenses in the appendix at camera-ready.]

\item {\bf New assets} \\
Question: Are new assets introduced in the paper well documented? \\
Answer: \answerYes{} \\
Justification: The instrument (code), the synthetic-Markov calibration corpora, and the toy checkpoints are documented in the repository README and the Reproducibility appendix.

\item {\bf Crowdsourcing and research with human subjects} \\
Question: N/A. \\
Answer: \answerNA{} \\
Justification: No human subjects or crowdsourcing.

\item {\bf Institutional review board (IRB) approvals} \\
Answer: \answerNA{} \\
Justification: No human-subjects research.

\item {\bf Declaration of LLM usage} \\
Question: Does the paper describe LLM usage if it was an important, original, or non-standard component? \\
Answer: \answerYes{} \\
Justification: Pretrained LLMs are the object of study (the CA rule), documented throughout. [TODO: disclose any LLM assistance in writing/experimentation per venue policy.]
\end{enumerate}
